% Compact 20-page dissertation (excl. appendixes and illustrations)
\documentclass[a4paper,oneside,11pt]{report}

\input{template-preamble}

\usepackage{newpxtext,newpxmath}
\usepackage{inconsolata}
\usepackage{microtype}
\usepackage{float}

\dissertationtitle{Evidence-Weighted and Opponent-Shaping\\
  Policy Gradients in Stochastic Games}

\candnumber{007}

\parskip1ex

\begin{document}

\dissertationtitlepage

\begin{dissertationsummary}
We study the convergence of policy gradient methods to Nash equilibria in multi-agent stochastic games, building on Giannou et al.\ (2022). We introduce two modifications that provably improve convergence: (1)~an \emph{evidence-weighted} variant where agents scale gradient updates by a confidence weight, reducing the effective variance by a factor of $\mathrm{HM}(V)/\mathrm{AM}(V) \leq 1$; and (2)~an \emph{opponent-shaping} variant incorporating Foerster et al.'s LOLA mechanism, which, under a spectral reinforcement condition, strictly enlarges the basin of attraction. The two improvements are orthogonal and compose in a combined algorithm. Both results are formalised in Lean~4 and validated empirically on matrix games.
\end{dissertationsummary}

\dissertationtableofcontents

% ======================================================================
% CHAPTER 1: INTRODUCTION
% ======================================================================
\chapter{Introduction}
\label{ch:introduction}

\section{Motivation and Contributions}

Learning in multi-agent stochastic games is fundamentally harder than single-agent optimisation: each agent's reward depends on all agents' actions, which are simultaneously changing as those agents learn. The convergence properties of policy gradient methods remain poorly understood outside specific game classes~\cite{giannou2022}.

Giannou et al.~\cite{giannou2022} established local convergence of policy gradient methods to Nash equilibria in \emph{general} stochastic games, with an $\mathcal{O}(1/\sqrt{n})$ rate for second-order stationary (SOS) policies under REINFORCE. This dissertation asks: \emph{can we do better?}

We introduce two modifications and prove each yields a measurable improvement:

\begin{enumerate}
    \item \textbf{Evidence-weighted policy gradient (EWPG).} Agents scale gradient updates by a confidence measure $V_i$ reflecting data quality. Under a heterogeneous variance model $\sigma_i^2 \propto 1/V_i$, the effective variance improves by
    $\operatorname{HM}(V)/\operatorname{AM}(V) \leq 1$.
    The rate exponent $\mathcal{O}(1/\sqrt{n})$ is preserved; the improvement is in the leading constant.

    \item \textbf{Opponent-shaping policy gradient (LOLA-PG).} Building on LOLA~\cite{foerster2018lola}, agents differentiate through opponents' anticipated learning steps. We prove that: (a)~with annealed shaping strength $\lambda_n \to 0$, convergence guarantees are preserved; (b)~under a spectral reinforcement condition, the SOS parameter increases from $\mu$ to $\mu + \lambda\mu_H$, \emph{strictly enlarging} the basin of attraction. This is, to our knowledge, the first convergence guarantee for LOLA-type methods in general stochastic games.
\end{enumerate}

The two improvements are orthogonal: evidence weighting reduces the variance prefactor, opponent shaping increases the contraction rate.

\section{Related Work}

\paragraph{Policy gradient convergence.} Single-agent convergence is well understood~\cite{agarwal2021theory,mei2020global}. In multi-agent settings, convergence has been established in zero-sum games~\cite{daskalakis2018training}, potential games~\cite{leonardos2022global}, and Markov potential games~\cite{cen2021fast}. Giannou et al.~\cite{giannou2022} extended guarantees to general stochastic games by focusing on equilibrium policies satisfying a second-order stationarity condition.

\paragraph{Opponent shaping.} LOLA~\cite{foerster2018lola} differentiates through opponents' anticipated updates, achieving cooperation in the Iterated Prisoner's Dilemma. Three lines of follow-up work lack general-game convergence guarantees: Stable Opponent Shaping~\cite{letcher2019} converges to stable fixed points but without rates or basin analysis; COLA~\cite{willi2022cola} addresses the inconsistency of simultaneous shaping but restricts to differentiable games; model-free opponent shaping~\cite{lu2022model} eliminates Hessian computation but provides no convergence theory. None provides last-iterate Nash convergence in general stochastic games.

\paragraph{Meta-learning in MARL.} Al-Shedivat et al.~\cite{alshedivat2018} framed multi-agent interaction as meta-learning. Kim et al.~\cite{kim2021policygradientalgorithmlearning} derived the Meta-MAPG gradient decomposition (Appendix~B), showing LOLA is a first-order approximation of the full three-term meta-gradient.

\paragraph{Positioning.} Giannou et al.\ provide the convergence framework but assume homogeneous agents and no opponent modelling. Foerster et al.\ provide opponent shaping but no convergence guarantees. We fill both gaps: EWPG handles heterogeneous evidence quality, LOLA-PG provides the first convergence guarantee for opponent-shaping methods in general stochastic games.

\section{Outline}

Chapter~\ref{ch:framework} develops the learning framework: the embedding-predictor decomposition, gradient structure, and risk decomposition that motivates the convergence results. Chapter~\ref{ch:results} presents the two main results with full proofs. Chapter~\ref{ch:simulations} validates them experimentally. Chapter~\ref{ch:conclusion} concludes. Appendices contain the Policy Gradient Theorem proof and Meta-MAPG decomposition. Both main results are formalised in Lean~4.


\chapter{The Learning Framework}
\label{ch:framework}

This chapter develops a mathematical framework for multi-agent learning that motivates the convergence results of Chapter~\ref{ch:results}. We define the data pipeline (\S\ref{sec:pipeline}), derive the gradient decomposition into embedding and predictor components (\S\ref{sec:gradient-decomp}), decompose the expected risk over binary partitions (\S\ref{sec:risk-decomp}), and extend the framework to multiple agents (\S\ref{sec:multi-agent-framework}), arriving at a taxonomy that classifies different learning algorithms and ensemble methods by their comparative advantage.


%% ========================================================================
\section{The Data Pipeline}
\label{sec:pipeline}
%% ========================================================================

We model a learner as a composition of two parametric maps operating on structured data.

\begin{definition}[Data pipeline]
\label{def:pipeline}
A learning pipeline consists of:
\begin{enumerate}[label=(\roman*)]
    \item An \emph{observation space} $O$ and an observation $\mathcal{X} \in O$;
    \item A \emph{proposition map} $\tau: O \to \Pi$, producing a structured representation $X = \tau(\mathcal{X})$ in a proposition space $\Pi$;
    \item A parametric \emph{embedding} $\phi_\psi: \Pi \to \mathbb{R}^m$ with parameters $\psi \in \mathbb{R}^p$, producing a feature vector $\mathbf{X} = \phi_\psi(X)$;
    \item A \emph{true causal map} $f: \mathbb{R}^m \to \mathbb{R}^k$ (unknown), defining the target $Y = f(\mathbf{X})$;
    \item A parametric \emph{predictor} $\hat{f}_\theta: \mathbb{R}^m \to \mathbb{R}^k$ with parameters $\theta \in \mathbb{R}^l$, producing predictions $\hat{Y} = \hat{f}_\theta(\mathbf{X})$;
    \item A loss function $\mathcal{L}: \mathbb{R}^k \times \mathbb{R}^k \to \mathbb{R}$, continuously differentiable, with $\mathcal{L}(Y, \hat{Y}) = \mathcal{L}(f(\mathbf{X}_\psi), \hat{f}_\theta(\mathbf{X}_\psi))$.
\end{enumerate}
We write $\beta = (\theta, \psi) \in \mathbb{R}^{l+p}$ for the full parameter vector.
\end{definition}

The pipeline separates two distinct learning tasks: \emph{what to represent} (the embedding $\phi_\psi$) and \emph{how to predict} (the predictor $\hat{f}_\theta$). This separation is natural in modern ML architectures: in a transformer, $\psi$ are the embedding/encoder parameters and $\theta$ are the prediction head parameters; in a GBM with engineered features, $\psi$ encodes the feature construction and $\theta$ the tree ensemble.


%% ========================================================================
\section{Gradient Decomposition}
\label{sec:gradient-decomp}
%% ========================================================================

The embedding-predictor separation produces structurally different gradients for each parameter group.

\begin{proposition}[Predictor gradient]
\label{prop:theta-grad}
Since $Y = f(\mathbf{X}_\psi)$ does not depend on $\theta$:
\begin{equation}
\label{eq:theta-grad}
    \nabla_\theta \mathcal{L} = \frac{\partial \mathcal{L}}{\partial \hat{Y}} \cdot \frac{\partial \hat{f}_\theta(\mathbf{X}_\psi)}{\partial \theta} \in \mathbb{R}^{1 \times l}.
\end{equation}
This is standard backpropagation: only the learned predictor contributes.
\end{proposition}

\begin{proposition}[Embedding gradient]
\label{prop:psi-grad}
Both arguments of $\mathcal{L}$ depend on $\psi$ through $\mathbf{X}_\psi = \phi_\psi(X)$. By the multivariate chain rule:
\begin{equation}
\label{eq:psi-grad}
\boxed{
    \nabla_\psi \mathcal{L} = \underbrace{\frac{\partial \mathcal{L}}{\partial Y} \cdot \frac{\partial f}{\partial \mathbf{X}} \cdot \frac{\partial \phi_\psi}{\partial \psi}}_{\text{through } f \text{ (inaccessible)}} + \underbrace{\frac{\partial \mathcal{L}}{\partial \hat{Y}} \cdot \frac{\partial \hat{f}_\theta}{\partial \mathbf{X}} \cdot \frac{\partial \phi_\psi}{\partial \psi}}_{\text{through } \hat{f}_\theta \text{ (computable)}} \in \mathbb{R}^{1 \times p}.
}
\end{equation}
\end{proposition}

\begin{proof}
$\mathcal{L}(\psi) = \mathcal{L}(f(\phi_\psi(X)),\, \hat{f}_\theta(\phi_\psi(X)))$. Both $f \circ \phi_\psi$ and $\hat{f}_\theta \circ \phi_\psi$ depend on $\psi$:
\[
    \nabla_\psi \mathcal{L} = \frac{\partial \mathcal{L}}{\partial Y} \nabla_\psi Y + \frac{\partial \mathcal{L}}{\partial \hat{Y}} \nabla_\psi \hat{Y} = \frac{\partial \mathcal{L}}{\partial Y} \frac{\partial f}{\partial \mathbf{X}} \frac{\partial \phi_\psi}{\partial \psi} + \frac{\partial \mathcal{L}}{\partial \hat{Y}} \frac{\partial \hat{f}_\theta}{\partial \mathbf{X}} \frac{\partial \phi_\psi}{\partial \psi}. \qedhere
\]
\end{proof}

\begin{remark}[The inaccessible term]
\label{rem:inaccessible}
The first term in~\eqref{eq:psi-grad} requires $\partial f / \partial \mathbf{X}$---the Jacobian of the \emph{true} causal map. If $f$ is unknown (the standard case), this term is inaccessible. In practice, the embedding gradient is approximated by the second term alone, introducing a systematic bias. This bias vanishes only when $\hat{f}_\theta = f$ (perfect prediction) or when $\partial \mathcal{L}/\partial Y = 0$ (which requires special loss structure). The inaccessible term can be estimated if the causal structure of $f$ is known---e.g., through Pearl's do-calculus when $f$ decomposes along a known DAG.
\end{remark}


%% ========================================================================
\section{Risk Decomposition}
\label{sec:risk-decomp}
%% ========================================================================

\subsection{Expected risk and the exploration-exploitation decomposition}

When observations $\mathcal{X}$ are drawn from a distribution $p_\theta$ over $O$ that depends on $\theta$ (e.g., in RL, where the policy determines which states are visited), the expected risk is:
\begin{equation}
\label{eq:expected-risk}
    R(\theta, \psi) = \mathbb{E}_{\mathcal{X} \sim p_\theta}\!\left[\mathcal{L}\bigl(f(\phi_\psi(\tau(\mathcal{X}))),\; \hat{f}_\theta(\phi_\psi(\tau(\mathcal{X})))\bigr)\right].
\end{equation}

\begin{proposition}[Exploration-exploitation decomposition]
\label{prop:reinforce-decomp}
If $p_\theta$ depends on $\theta$, the gradient decomposes as:
\begin{equation}
\label{eq:reinforce-decomp}
\boxed{
    \nabla_\theta R = \underbrace{\mathbb{E}_{p_\theta}\!\left[\mathcal{L}(Y, \hat{Y}) \cdot \nabla_\theta \log p_\theta(\mathcal{X})\right]}_{\text{exploration (REINFORCE)}} + \underbrace{\mathbb{E}_{p_\theta}\!\left[\nabla_\theta \mathcal{L}(Y, \hat{Y})\right]}_{\text{exploitation (backprop)}}.
}
\end{equation}
\end{proposition}

\begin{proof}
$\nabla_\theta R = \nabla_\theta \int_O \mathcal{L} \cdot p_\theta \, d\mathcal{X} = \int_O \mathcal{L} \cdot \nabla_\theta p_\theta \, d\mathcal{X} + \int_O \nabla_\theta \mathcal{L} \cdot p_\theta \, d\mathcal{X}$.
Using the log-derivative trick $\nabla_\theta p_\theta = p_\theta \cdot \nabla_\theta \log p_\theta$ on the first term gives~\eqref{eq:reinforce-decomp}. \qedhere
\end{proof}

The exploration term adjusts \emph{where} in $O$ the learner looks (score-weighted loss); the exploitation term adjusts \emph{how well} the learner predicts given what it sees. This is precisely the Policy Gradient Theorem (Appendix~\ref{app:pgt}) when $\theta$ parametrises a policy and $\mathcal{L}$ is the negative reward.


\subsection{Binary partition of the risk}

Let $\Pi = W \cup S$ be any binary partition of the proposition space (e.g., true vs.\ false propositions, in-distribution vs.\ out-of-distribution, easy vs.\ hard examples). The expected risk decomposes:

\begin{equation}
\label{eq:ws-decomp}
    R(\theta, \psi) = \underbrace{P_W \cdot \mathbb{E}[\mathcal{L} \mid X \in W]}_{R_W} + \underbrace{P_S \cdot \mathbb{E}[\mathcal{L} \mid X \in S]}_{R_S},
\end{equation}
where $P_W = P_\theta(\tau^{-1}(W))$ and $P_S = 1 - P_W$.

\begin{proposition}[Risk Taylor expansion over partition]
\label{prop:taylor-partition}
Let $Y^* \in \mathbb{R}^k$ be a reference prediction and $\delta = \hat{Y} - Y^*$. Taylor-expanding $\mathcal{L}(Y, Y^* + \delta)$ in $\delta$:
\begin{equation}
\label{eq:risk-taylor}
\boxed{
\begin{aligned}
    R(\theta, \psi) &\approx \underbrace{P_W \mathbb{E}[\mathcal{L}(Y, Y^*) \mid W] + P_S \mathbb{E}[\mathcal{L}(Y, Y^*) \mid S]}_{\text{irreducible loss at reference}} \\
    &\quad + \underbrace{P_W \mathbb{E}[\nabla_{\hat{Y}} \mathcal{L}^\top \delta \mid W] + P_S \mathbb{E}[\nabla_{\hat{Y}} \mathcal{L}^\top \delta \mid S]}_{\text{bias (first-order correction)}} \\
    &\quad + \underbrace{\tfrac{1}{2} P_W \mathbb{E}[\delta^\top H_{\hat{Y}} \delta \mid W] + \tfrac{1}{2} P_S \mathbb{E}[\delta^\top H_{\hat{Y}} \delta \mid S]}_{\text{variance / curvature (second-order correction)}},
\end{aligned}
}
\end{equation}
where $H_{\hat{Y}} = \nabla^2_{\hat{Y}} \mathcal{L}(Y, Y^*)$ is the Hessian of the loss.
\end{proposition}

Each row decomposes separately over the partition. This gives six terms, each with a distinct interpretation: three sources of error (irreducible, bias, variance) $\times$ two subpopulations ($W, S$).


%% ========================================================================
\section{Extension to Multiple Agents}
\label{sec:multi-agent-framework}
%% ========================================================================

Now let there be $N$ agents, indexed $i = 1, \ldots, N$. Each agent $i$ has its own:
\begin{itemize}
    \item Embedding $\phi_{\psi_i}: \Pi \to \mathbb{R}^{m_i}$ with parameters $\psi_i$,
    \item Predictor $\hat{f}_{\theta_i}$ with parameters $\theta_i$,
    \item Sampling distribution $p_{\theta_i}$ over $O$,
    \item \emph{Weight of evidence} $V_i(\mathcal{X}) > 0$, measuring how much data supports its predictions.
\end{itemize}

\subsection{Evidence-weighted multi-agent risk}

Define the \emph{pooled weight of evidence}:
\begin{equation}
\label{eq:pooled-evidence}
    V_{\mathrm{pool}}(\mathcal{X}) = \sum_{i=1}^N \alpha_i V_i(\mathcal{X}),
\end{equation}
where $\alpha_i \geq 0$ are credibility weights. The multi-agent risk augments the prediction loss with an evidence-seeking penalty:
\begin{equation}
\label{eq:multi-agent-risk}
    R_{\mathrm{multi}}(\Theta, \Psi) = \sum_{i=1}^N \omega_i R_i(\theta_i, \psi_i) + \mu \cdot V_{\mathrm{pool}}^{-1},
\end{equation}
where $\omega_i$ are agent importance weights and $\mu > 0$ trades off prediction quality against evidentiary robustness. The gradient for agent $i$ decomposes as:
\begin{equation}
\label{eq:five-component}
\boxed{
    \nabla_{\theta_i} R_{\mathrm{multi}} = \omega_i\Bigl[\underbrace{\mathbb{E}[\mathcal{L}_i \cdot s_{\theta_i}]}_{\text{exploration}} + \underbrace{\mathbb{E}[\nabla_{\theta_i} \mathcal{L}_i]}_{\text{exploitation}}\Bigr] + \underbrace{\omega_i \mu(-V_i^{-2} \nabla_{\theta_i} V_i)}_{\text{evidence seeking}} + \underbrace{\sum_{j \neq i} \frac{\partial R_i}{\partial \theta_j} \frac{\partial \theta_j'}{\partial \theta_i}}_{\text{opponent shaping}},
}
\end{equation}
where $s_{\theta_i} = \nabla_{\theta_i} \log p_{\theta_i}$ is the score function.


\subsection{Risk decomposition for multi-agent systems}
\label{subsec:six-way}

Define agent $i$'s \emph{effective domain} as the set of propositions for which it has sufficient evidence: $\Pi_N^i = \{X \in \Pi \mid V_i(X) \geq V_{\min}\}$, and let $\Pi_U^i = \Pi \setminus \Pi_N^i$ be the complement. Define the \emph{collective blind spot} as $B = \bigcap_i B_i$, where $B_i$ is the set of propositions agent $i$ cannot resolve.

The multi-agent risk admits a six-way decomposition:
\begin{equation}
\label{eq:six-way}
\boxed{
    R_{\mathrm{multi}} = \underbrace{R_{W \cap \Pi_N \cap B^c} + R_{S \cap \Pi_N \cap B^c}}_{\text{learnable (gradient descent)}} + \underbrace{R_{W \cap \Pi_N \cap B} + R_{S \cap \Pi_N \cap B}}_{\text{blind-spot-limited}} + \underbrace{R_{W \cap \Pi_U} + R_{S \cap \Pi_U}}_{\text{evidence-limited}}.
}
\end{equation}

Only the first two terms are reducible by standard gradient descent. The middle pair requires agents to communicate complementary information---each agent $i$ may resolve propositions in $B_j \setminus B_i$ that agent $j$ cannot, and vice versa. The last pair requires increasing the weight of evidence through more data or better observation.

\begin{remark}[Connection to the convergence results]
\label{rem:connection}
The two contributions of Chapter~\ref{ch:results} map precisely onto this decomposition:
\begin{itemize}
    \item \textbf{Evidence-weighted PG} addresses the evidence-limited terms ($R_{W \cap \Pi_U}$, $R_{S \cap \Pi_U}$): agents with heterogeneous evidence quality scale their gradient updates to avoid acting on unreliable estimates, reducing variance in the Lyapunov analysis by $\operatorname{HM}(V)/\operatorname{AM}(V)$.
    \item \textbf{Opponent-shaping PG} addresses the blind-spot-limited terms ($R_{W \cap \Pi_N \cap B}$, $R_{S \cap \Pi_N \cap B}$): by modelling opponents' learning dynamics, agents expand their effective convergence basin beyond what independent learning achieves.
\end{itemize}
The learnable terms are already addressed by standard PG~\cite{giannou2022}. A complete multi-agent gradient would additionally include a consensus alignment term for the evidence-limited terms, which we leave to future work.
\end{remark}


%% ========================================================================
\section{Instantiation in Stochastic Games}
\label{sec:instantiation}
%% ========================================================================

We now show how the general learning framework specialises to the episodic stochastic game setting of Giannou et al.~\cite{giannou2022}, which is the concrete domain for the convergence results of Chapter~\ref{ch:results}.

\paragraph{Mapping the pipeline.} In an $N$-player stochastic game $\mathcal{G} = (\mathcal{S}, \mathcal{N}, \{\mathcal{A}_i\}, P, \zeta, \rho)$:
\begin{center}
\renewcommand{\arraystretch}{1.3}
\begin{tabular}{lll}
\hline
\textbf{Framework} & \textbf{Stochastic game} & \textbf{Role} \\
\hline
Observation $\mathcal{X}$ & Episode trajectory $\tau = (s_t, \alpha_t, r_t)_{t \leq T}$ & Raw experience \\
Proposition $X = \tau(\mathcal{X})$ & State-action-reward sequence & Structured data \\
Embedding $\phi_\psi(X)$ & State encoding (identity in tabular case) & Feature representation \\
Predictor $\hat{f}_\theta$ & Policy $\pi_\theta: \mathcal{S} \to \Delta(\mathcal{A})$ & Decision rule \\
True map $f$ & Optimal policy $\pi^*$ (Nash) & Unknown target \\
Loss $\mathcal{L}$ & Negative reward $-R_i(s, \alpha_i, \alpha_{-i})$ & Per-step cost \\
Sampling dist.\ $p_\theta$ & Trajectory distribution under $\pi_\theta$ & Exploration \\
Evidence weight $V_i$ & Sample budget or inverse gradient variance & Data quality \\
\hline
\end{tabular}
\end{center}

\paragraph{Recovering the Policy Gradient Theorem.} Under this mapping, the exploration-exploitation decomposition~\eqref{eq:reinforce-decomp} becomes:
\begin{equation}
\label{eq:pgt-from-framework}
    \nabla_\theta J(\theta) = \mathbb{E}_{\tau \sim \pi_\theta}\!\left[R(\tau) \cdot \nabla_\theta \log \pi_\theta(\tau)\right],
\end{equation}
which is the standard Policy Gradient Theorem (proved in Appendix~\ref{app:pgt}). The REINFORCE estimator samples a single trajectory $\tau$ and uses $R(\tau) \cdot \sum_t \nabla_\theta \log \pi_\theta(a_t | s_t)$ as a stochastic gradient. The high variance of this estimator---which scales as $\mathcal{O}(1/\kappa)$ where $\kappa = \min_{s,a} \pi(a|s)$ is the minimum action probability---is exactly the noise term $U_n$ in the Giannou et al.\ framework.

\paragraph{The episodic random-stopping formulation.} Giannou et al.\ consider games where episodes terminate with probability $\zeta_{s,\alpha} > 0$ at each step. The value function for agent $i$ is $V_{i,s}(\pi) = \mathbb{E}_{\tau \sim \text{MDP}}[\sum_{t=0}^{T(\tau)} R_i(s_t, \alpha_t) \mid s_0 = s]$, and the expected risk under the initial state distribution $\rho$ is $V_{i,\rho}(\pi) = \mathbb{E}_{s \sim \rho}[V_{i,s}(\pi)]$. This is the multi-agent risk $R_{\mathrm{multi}}$ of~\eqref{eq:multi-agent-risk} with the evidence-seeking and alignment terms set to zero.

\paragraph{The inaccessible gradient in RL.} The inaccessible term $\partial f / \partial \mathbf{X}$ in the embedding gradient~\eqref{eq:psi-grad} corresponds to the derivative of the \emph{optimal} policy with respect to state features---information that requires knowing the optimal policy itself. In practice, this term is estimated through temporal-difference learning or ignored entirely. The multi-agent setting adds a further inaccessible term: the derivative of agent $i$'s reward through agent $j$'s policy, which LOLA approximates via first-order Taylor expansion (see \S\ref{sec:lola-conv}).


\input{chapters/ch03-convergence}


\chapter{Experiments}
\label{ch:simulations}

This chapter validates the theoretical results of Chapter~\ref{ch:results} through six experiments of increasing depth. We verify the AM-HM variance improvement (\S\ref{sec:exp-variance}), the convergence rate and constant (\S\ref{sec:exp-rate}), the spectral reinforcement condition (\S\ref{sec:exp-spectral}), the basin of attraction enlargement (\S\ref{sec:exp-basin}), the combined EW-LOLA-PG algorithm (\S\ref{sec:exp-combined}), and the robustness across game types (\S\ref{sec:exp-multigame}).

All experiments use the REINFORCE gradient estimator (Algorithm~1 of~\cite{giannou2022}) with step size $\gamma_n = \gamma/(n+m)^p$, $\gamma = 0.5$, $m = 10$, $p = 0.75$, and simplex projection. Source code is at \texttt{simulations/}.


%% ========================================================================
\section{Gradient Variance Measurement}
\label{sec:exp-variance}
%% ========================================================================

We first test Theorem~\ref{thm:am-hm-ch3} by directly measuring gradient variance.

\paragraph{Setup.} Fix policies near the Nash equilibrium of Matching Pennies ($\pi_1 = (0.52, 0.48)$, $\pi_2 = (0.48, 0.52)$). For each evidence ratio $r = V_1/V_2 \in \{1, 2, 3, 5, 8, 10, 15, 20, 50\}$, draw 2000 REINFORCE gradient samples with additive noise $\sigma_i = 2/\sqrt{V_i}$, and compute the trace of the gradient covariance matrix for both standard PG and EWPG.

\begin{figure}[H]
\centering
\includegraphics[width=0.48\textwidth]{figures/extended/gradient_variance.png}
\caption{\emph{Left}: Absolute gradient variance for standard PG (red) and EWPG (green) as a function of evidence ratio. Standard PG variance is dominated by the low-evidence agent; EWPG downweights it. \emph{Right}: Empirical variance ratio (blue) vs.\ theoretical HM/AM prediction (black). Close agreement across two orders of magnitude in $r$.}
\label{fig:gradient-variance}
\end{figure}

The empirical ratio tracks the theoretical curve $4r/(1+r)^2$ closely (Figure~\ref{fig:gradient-variance}, right), confirming that the variance improvement is not an artefact of the convergence dynamics but a direct consequence of the weighting scheme applied to the gradient estimator.


%% ========================================================================
\section{Convergence Rate Verification}
\label{sec:exp-rate}
%% ========================================================================

Theorem~\ref{thm:ewpg-conv-ch3} predicts that EWPG preserves the $\mathcal{O}(1/\sqrt{n})$ rate exponent while improving the leading constant by $\operatorname{HM}(V)/\operatorname{AM}(V)$.

\paragraph{Setup.} Matching Pennies, 10{,}000 episodes, 100 runs, $\gamma_n = 0.5/(n+10)^{0.75}$. Three configurations: Standard PG ($r=1$), EWPG with $r=10$, EWPG with $r=100$.

\begin{figure}[H]
\centering
\includegraphics[width=0.95\textwidth]{figures/extended/convergence_rate.png}
\caption{\emph{Left}: Log-log plot of $\mathbb{E}[\|\pi_n - \pi^*\|^2]$ vs.\ episode number. All three methods exhibit the same slope ($\approx -0.5$, dashed lines), confirming the rate exponent is preserved. The vertical offset decreases with increasing $r$, confirming the constant improvement. \emph{Right}: Empirical constant ratio $C_w/C_{\mathrm{std}}$ vs.\ theoretical HM/AM across evidence ratios up to $r = 100$.}
\label{fig:convergence-rate}
\end{figure}

We verify the rate exponent by fitting a line on the log-log plot (last 60\% of data). All three configurations yield slopes close to $-0.5$, consistent with the $\mathcal{O}(1/\sqrt{n})$ prediction. The vertical offset between curves quantifies the constant improvement, which matches the theoretical HM/AM ratio (Figure~\ref{fig:convergence-rate}, right).


%% ========================================================================
\section{Spectral Reinforcement Analysis}
\label{sec:exp-spectral}
%% ========================================================================

Theorem~\ref{thm:basin-ch3} requires the spectral reinforcement condition ($\mu_H > 0$) on the opponent-shaping Hessian $H$. We compute $H$ numerically for three games and verify when the condition holds.

\paragraph{Setup.} For each 2$\times$2 game, compute the opponent-shaping Hessian $H(\pi^*)$ at the interior Nash equilibrium via finite differences. Extract the symmetric part $S_H = (H + H^\top)/2$ and its eigenvalues.

\begin{figure}[H]
\centering
\includegraphics[width=0.95\textwidth]{figures/extended/spectral_reinforcement.png}
\caption{Effective SOS parameter $\mu_{\mathrm{LOLA}} = \mu + \lambda\mu_H$ vs.\ LOLA strength $\lambda$ for three games. Green region: $\mu_{\mathrm{LOLA}} > 0$ (convergence guaranteed). Matching Pennies (zero-sum): strong spectral reinforcement ($\mu_H > 0$), $\mu_{\mathrm{LOLA}}$ increases with $\lambda$. Battle of the Sexes and Stag Hunt: weaker or absent reinforcement.}
\label{fig:spectral}
\end{figure}

Matching Pennies exhibits strong spectral reinforcement: the opponent-shaping Hessian has $\mu_H > 0$, meaning LOLA monotonically improves the contraction rate as $\lambda$ increases (Figure~\ref{fig:spectral}, left). This confirms that zero-sum games are the natural domain for LOLA, as the gradient Jacobian has a large antisymmetric (rotational) component that LOLA converts to contraction.


%% ========================================================================
\section{Basin of Attraction}
\label{sec:exp-basin}
%% ========================================================================

\subsection{Basin mapping}

We map the convergence basin by sweeping initial conditions over a $25 \times 25$ grid in $[0.05, 0.95]^2$, running 8 trials per point with noise $\sigma = 0.3$.

\begin{figure}[H]
\centering
\includegraphics[width=0.48\textwidth]{figures/lola/basin_matching_pennies.png}
\includegraphics[width=0.48\textwidth]{figures/lola/basin_battle_of_the_sexes.png}
\caption{Basin of attraction: Standard PG (left), LOLA-PG (centre), difference (right). White star: Nash equilibrium. LOLA approximately doubles the basin in Matching Pennies.}
\label{fig:basin}
\end{figure}

\subsection{Basin size vs.\ LOLA strength}

We sweep constant $\lambda \in \{0, 0.1, 0.2, 0.5, 1.0, 1.5, 2.0\}$ and measure basin area.

\begin{figure}[H]
\centering
\includegraphics[width=0.48\textwidth]{figures/extended/annealing_schedules.png}
\caption{\emph{Left}: Convergence under different LOLA annealing schedules. Constant $\lambda$ (purple) converges fastest initially but has residual bias. Annealed schedules $\lambda_n \propto 1/n^r$ converge to zero distance, with $r = 0.75$ matching the step-size exponent (Theorem~\ref{thm:annealed-ch3}). \emph{Right}: Basin size (fraction of state space converging) vs.\ constant $\lambda$. Basin grows monotonically, consistent with $\mu_{\mathrm{LOLA}} = \mu + \lambda\mu_H$.}
\label{fig:annealing}
\end{figure}

\subsection{Trajectory visualisation}

\begin{figure}[H]
\centering
\includegraphics[width=0.48\textwidth]{figures/lola/trajectories_matching_pennies.png}
\hfill
\includegraphics[width=0.48\textwidth]{figures/lola/trajectories_battle_of_the_sexes.png}
\caption{Sample trajectories (green = converged, red = diverged). Under standard PG, Matching Pennies trajectories orbit the Nash equilibrium. LOLA converts these orbits into convergent spirals.}
\label{fig:trajectories}
\end{figure}

The trajectory plots (Figure~\ref{fig:trajectories}) reveal the mechanism: in Matching Pennies, the gradient field has a strong rotational component at Nash. Standard PG orbits. LOLA adds a radial (contractive) component via the opponent-shaping term, converting circular motion to an inward spiral. This is the geometric content of the spectral reinforcement theorem.


%% ========================================================================
\section{Combined EW-LOLA-PG}
\label{sec:exp-combined}
%% ========================================================================

Theorem~\ref{thm:combined-ch3} predicts that evidence weighting and opponent shaping compose orthogonally. We test all four combinations.

\paragraph{Setup.} Matching Pennies, $V = (10, 1)$, $\lambda = 0.5$ (constant), 5{,}000 episodes, 80 runs.

\begin{figure}[H]
\centering
\includegraphics[width=0.95\textwidth]{figures/extended/combined_ewlola.png}
\caption{\emph{Left}: Convergence curves for all four methods. EW-LOLA-PG (gold) dominates, inheriting lower variance from EWPG and faster contraction from LOLA. \emph{Right}: Mean distance in last 500 episodes. Ratios relative to Standard PG shown above bars. EW-LOLA-PG achieves the best final performance, confirming orthogonal composition.}
\label{fig:combined}
\end{figure}

The combined algorithm achieves the lowest asymptotic distance (Figure~\ref{fig:combined}), and the improvement factors are approximately multiplicative: if EWPG gives factor $a$ and LOLA gives factor $b$, EW-LOLA gives approximately $a \times b$. This confirms that the two mechanisms act on independent terms of the Lyapunov recursion (variance and contraction respectively).


%% ========================================================================
\section{Multi-Game Robustness}
\label{sec:exp-multigame}
%% ========================================================================

We test EWPG across three structurally different games to assess robustness.

\begin{figure}[H]
\centering
\includegraphics[width=0.95\textwidth]{figures/extended/multigame_comparison.png}
\caption{EWPG (green) vs.\ Standard PG (red) on Matching Pennies (zero-sum), Rock-Paper-Scissors (zero-sum, $3 \times 3$), and Stag Hunt (coordination). EWPG consistently maintains lower variance. The improvement is most pronounced in zero-sum games where the gradient field has large rotational components.}
\label{fig:multigame}
\end{figure}

EWPG improves convergence quality consistently across all tested games (Figure~\ref{fig:multigame}). The improvement is largest in zero-sum games (Matching Pennies, RPS) where gradient variance is the dominant obstacle, and present but smaller in the coordination game (Stag Hunt).


%% ========================================================================
\section{Discussion}
\label{sec:exp-discussion}
%% ========================================================================

The experiments confirm all theoretical predictions:

\begin{enumerate}
    \item \textbf{Variance improvement} (Theorem~\ref{thm:am-hm-ch3}): Empirical variance ratio matches HM/AM across $r \in [1, 50]$. Direct gradient variance measurement and convergence constant measurement agree.
    \item \textbf{Rate preservation} (Theorem~\ref{thm:ewpg-conv-ch3}): Log-log slopes $\approx -0.5$ for all configurations. Only the constant changes.
    \item \textbf{Spectral reinforcement} (Definition~\ref{def:spectral-ch3}): Zero-sum games have $\mu_H > 0$; coordination games do not. Matches Proposition~\ref{prop:sufficient-ch3}.
    \item \textbf{Basin enlargement} (Theorem~\ref{thm:basin-ch3}): Basin size grows monotonically with $\lambda$ in Matching Pennies. Trajectory plots confirm rotation-to-contraction conversion.
    \item \textbf{Annealing} (Theorem~\ref{thm:annealed-ch3}): Annealed LOLA converges to exact Nash; constant LOLA has residual bias $O(\lambda)$, as predicted by Theorem 2 of the LOLA analysis.
    \item \textbf{Orthogonal composition} (Theorem~\ref{thm:combined-ch3}): EW-LOLA-PG improvement is approximately multiplicative.
\end{enumerate}

%% ========================================================================
\section{Real REINFORCE Variance}
\label{sec:exp-real-reinforce}
%% ========================================================================

A key critique of the preceding experiments is that variance is injected artificially via additive Gaussian noise scaled as $\sigma_i = C/\sqrt{V_i}$. This section eliminates the artificial noise and uses \emph{only} the variance that arises naturally from the REINFORCE estimator.

\paragraph{Setup.} We model heterogeneous evidence quality through \emph{sample budgets}: agent~1 draws $K_1$ trajectory samples per episode, agent~2 draws $K_2$. The REINFORCE variance scales as $\mathcal{O}(1/K_i)$ naturally (by the CLT), providing a non-artificial realisation of $\sigma_i^2 \propto 1/V_i$ with $V_i = K_i$. No additional noise is injected.

\begin{figure}[H]
\centering
\includegraphics[width=0.95\textwidth]{figures/roadmap/real_reinforce.png}
\caption{\emph{Left}: Standard PG with varying sample budgets $(K_1, K_2)$. Convergence improves with more samples, but heterogeneous budgets $(10,1)$ are limited by the low-sample agent. \emph{Right}: EWPG vs Standard PG with real REINFORCE variance ($K_1=10, K_2=1$). EWPG downweights the noisy agent, producing lower-variance convergence. No artificial noise is injected---all variance comes from the REINFORCE estimator.}
\label{fig:real-reinforce}
\end{figure}

This confirms that the EWPG improvement is not an artefact of the noise injection protocol: it holds under the native variance structure of the REINFORCE estimator when agents have genuinely different data quality (sample budgets).


%% ========================================================================
\section{Multi-State Stochastic Game}
\label{sec:exp-stochastic}
%% ========================================================================

We test EWPG on a 2-player, 2-state, 2-action stochastic game---the first experiment that exercises the full stochastic game framework of~\cite{giannou2022} rather than the single-state special case.

\paragraph{Game design.} State~0 is a modified Matching Pennies (zero-sum); State~1 is a Prisoner's Dilemma variant. State transitions depend on joint actions, with termination probability $\zeta = 0.1$. Policies are state-dependent: $\pi_i = (\pi_i(s=0), \pi_i(s=1)) \in \Delta(\mathcal{A}_i)^2$. The approximate Nash equilibrium is found by a long reference run (20{,}000 episodes, 5 samples/episode).

\begin{figure}[H]
\centering
\includegraphics[width=0.95\textwidth]{figures/roadmap/stochastic_game.png}
\caption{\emph{Left}: EWPG vs Standard PG on a 2-state stochastic game with sample budgets $K = (8, 1)$. EWPG converges to a lower distance. \emph{Right}: Per-state convergence under EWPG. The zero-sum state (State~0) converges more slowly due to its rotational dynamics, while the PD state (State~1) converges faster. This illustrates the asymmetric convergence predicted by the risk decomposition~\eqref{eq:risk-taylor}.}
\label{fig:stochastic-game}
\end{figure}

The multi-state experiment confirms that EWPG improves convergence in stochastic games with state transitions---not just in the single-state (normal-form) special case.


%% ========================================================================
\section{Adaptive Evidence Weight Estimation}
\label{sec:exp-adaptive}
%% ========================================================================

A central critique of EWPG is that the evidence weights $V_i$ are \emph{assumed}, not \emph{derived}. Here we show that $V_i$ can be estimated online from the running variance of each agent's gradient history, eliminating the need for oracle knowledge.

\paragraph{Adaptive estimation.} Agent $i$ maintains a window of its last 50 gradient estimates. The estimated variance is $\hat{\sigma}_i^2 = \mathrm{Var}(\{g_{i,n-50}, \ldots, g_{i,n}\})$. The adaptive evidence weight is $\hat{V}_i = 1/\hat{\sigma}_i^2$, normalised to $\hat{w}_i = \hat{V}_i / \max_j \hat{V}_j$.

\begin{figure}[H]
\centering
\includegraphics[width=0.95\textwidth]{figures/roadmap/adaptive_evidence.png}
\caption{\emph{Left}: Convergence of Standard PG, EWPG with oracle weights, and EWPG with adaptive weights. Adaptive EWPG approaches oracle performance after an initial learning phase. \emph{Centre}: Estimated weights $w_1, w_2$ over time. The adaptive mechanism correctly identifies that agent~1 (10 samples) deserves higher weight than agent~2 (1 sample), converging to near-oracle values. \emph{Right}: Final distance comparison.}
\label{fig:adaptive-evidence}
\end{figure}

The adaptive estimation tracks the oracle weights within $\sim$200 episodes and produces convergence performance close to the oracle EWPG. This addresses the ``tautological variance model'' critique: the heterogeneous variance model $\sigma_i^2 = C/V_i$ need not be assumed a priori---it can be estimated from the gradient history, and the AM-HM improvement still holds.


%% ========================================================================
\section{Summary and Limitations}
\label{sec:exp-summary}
%% ========================================================================

The nine experiments validate all theoretical predictions and address the principal critiques:

\begin{enumerate}
    \item \textbf{Variance improvement}: confirmed with both artificial noise (\S\ref{sec:exp-variance}) and real REINFORCE variance (\S\ref{sec:exp-real-reinforce}).
    \item \textbf{Rate preservation}: log-log slopes $\approx -0.5$ with improved constant (\S\ref{sec:exp-rate}).
    \item \textbf{Multi-state games}: EWPG improvement persists in stochastic games with state transitions (\S\ref{sec:exp-stochastic}).
    \item \textbf{Adaptive estimation}: evidence weights can be learned online, removing the oracle assumption (\S\ref{sec:exp-adaptive}).
    \item \textbf{Spectral reinforcement}: verified numerically per game (\S\ref{sec:exp-spectral}).
    \item \textbf{Orthogonal composition}: confirmed empirically (\S\ref{sec:exp-combined}).
\end{enumerate}

\paragraph{Remaining limitations.}
\begin{itemize}
    \item The multi-state experiment uses a 2-state game; larger state spaces (e.g., grid worlds) would strengthen the empirical case.
    \item Policies remain tabular; extension to neural network parametrisations is needed for deep MARL applications.
    \item LOLA requires payoff cross-derivatives, estimated here from known matrices rather than learned.
    \item Adaptive weight estimation adds overhead and requires a burn-in period ($\sim$50 episodes).
\end{itemize}



% ======================================================================
% CHAPTER 5: CONCLUSION
% ======================================================================
\chapter{Conclusion}
\label{ch:conclusion}

We extended the convergence framework of Giannou et al.~\cite{giannou2022} with two orthogonal modifications: evidence weighting (variance improvement by $\operatorname{HM}(V)/\operatorname{AM}(V)$) and opponent shaping (basin enlargement by $\mu \to \mu + \lambda\mu_H$). Both are formalised in Lean~4 and validated experimentally on matrix games, a multi-state stochastic game, and with both artificial and real REINFORCE variance.

\paragraph{Limitations.} The heterogeneous variance model $\sigma_i^2 = C/V_i$ is motivated by the REINFORCE variance structure (\S\ref{sec:ewpg}, eq.~\ref{eq:reinforce-variance}) and validated empirically (\S\ref{sec:exp-real-reinforce}, \S\ref{sec:exp-adaptive}), but a fully rigorous derivation from the score-function estimator for general stochastic games remains open. The spectral reinforcement condition is domain-dependent (\S\ref{sec:exp-spectral}). Both improvements are local. The improvement is in the constant, not the rate exponent---though the constant is often the binding constraint in practice.

\paragraph{Open questions.}
\begin{enumerate}
    \item \emph{Model-free LOLA}: the opponent-shaping term requires payoff cross-derivatives. A critic-based approximation that estimates these from trajectory data would make EW-LOLA-PG fully model-free.
    \item \emph{Consensus alignment}: the five-component gradient~\eqref{eq:five-component} includes an alignment term $\lambda \nabla_{\theta_i} D$ that penalises divergence between agents' representations. Formalising its convergence properties would complete the framework.
    \item \emph{Beyond SOS policies}: all guarantees require the Nash equilibrium to be second-order stationary. Extending to non-isolated or non-SOS equilibria (e.g., via lazy projection as in Giannou et al.\ Theorem~3) is an important direction.
    \item \emph{Parametric policies}: extending from tabular to neural-network-parametrised policies would connect the theory to deep MARL practice.
\end{enumerate}


% ======================================================================
% APPENDICES
% ======================================================================
\appendix

% Chapter 4: Proof of the Policy Gradient Theorem

\chapter{Proof of the Policy Gradient Theorem}
\label{ch:pgt-proof}

We now prove Theorem~\ref{thm:pgt}. The argument, due to \cite{sutton2018rl} (see also \cite{williams1992} for an earlier derivation in the REINFORCE setting), uses only the product rule and the Bellman equation, yet yields a remarkable result: the gradient of the performance measure can be written without differentiating the on-policy state distribution. We first state the assumptions required for the proof, then present the argument in full detail.

\section{Assumptions}
\label{sec:pgt-assumptions}

The proof of the Policy Gradient Theorem rests on the following conditions.

\begin{enumerate}[label=(\roman*)]
  \item \emph{Finite MDP.}\; The state space~$\mathcal{S}$ and action space~$\mathcal{A}$ are finite, as in Definition~\ref{def:mdp}. (The result extends to compact or countable spaces under additional regularity; see Remark~\ref{rem:measure-returns}.)
  \item \emph{Differentiable policy.}\; The parametrised policy $\pi(a\mid s,\bm{\theta})$ is differentiable with respect to~$\bm{\theta}$ for all $(s,a) \in \mathcal{S}\times\mathcal{A}$ and all~$\bm{\theta}$.
  \item \emph{Fixed environment dynamics.}\; The transition function~$\mathcal{P}$ and reward function~$\mathcal{R}$ do not depend on the policy parameters~$\bm{\theta}$. This is the single-agent assumption that will be relaxed in the multi-agent setting.
  \item \emph{Bounded rewards.}\; There exists $R_{\max}<\infty$ such that $|\mathcal{R}(s,a,s')|\le R_{\max}$ for all $(s,a,s')$. Together with $\gamma<1$, this ensures that all value functions are bounded and that the infinite series in the proof converge absolutely.
\end{enumerate}

\section{The Proof}
\label{sec:pgt-proof-body}

\begin{proof}[Proof of Theorem~\ref{thm:pgt}]
We suppress the explicit dependence on~$\bm{\theta}$, writing $\pi(a\mid s)$ for $\pi(a\mid s,\bm{\theta})$. All gradients are with respect to~$\bm{\theta}$.

\medskip\noindent\textbf{Step 1 (Product rule).}\quad
Recall from~\eqref{eq:v-from-q} that $v_\pi(s) = \sum_a \pi(a\mid s)\,q_\pi(s,a)$. Since both $\pi$ and~$q_\pi$ depend on~$\bm{\theta}$, differentiating via the product rule gives:
\begin{equation}\label{eq:grad-v-step1}
  \nabla v_\pi(s) = \sum_a \bigl[\nabla\pi(a\mid s)\,q_\pi(s,a) + \pi(a\mid s)\,\nabla q_\pi(s,a)\bigr].
\end{equation}
The first term inside the sum captures the direct effect of changing $\bm{\theta}$ on the action probabilities; the second captures the indirect effect through changes in the action-value function.

\medskip\noindent\textbf{Step 2 (Differentiating the Bellman equation).}\quad
To handle $\nabla q_\pi(s,a)$, we use the Bellman equation for the action-value function~\eqref{eq:bellman-q}:
\[
  q_\pi(s,a) = \sum_{s'\in\mathcal{S}} \mathcal{P}(s'\mid s,a)\bigl[\mathcal{R}(s,a,s') + \gamma\, v_\pi(s')\bigr].
\]
By assumption~(iii), neither $\mathcal{P}$ nor $\mathcal{R}$ depends on~$\bm{\theta}$. The only $\bm{\theta}$-dependence on the right-hand side is through~$v_\pi(s')$. Differentiating:
\begin{equation}\label{eq:grad-q}
  \nabla q_\pi(s,a)
  = \sum_{s'\in\mathcal{S}} \mathcal{P}(s'\mid s,a)\bigl[\underbrace{\nabla\mathcal{R}(s,a,s')}_{=\,0} + \gamma\,\nabla v_\pi(s')\bigr]
  = \gamma\sum_{s'\in\mathcal{S}}\mathcal{P}(s'\mid s,a)\,\nabla v_\pi(s').
\end{equation}
This is the step where the fixed-environment assumption is used most directly.

\medskip\noindent\textbf{Step 3 (Substitution and recursion).}\quad
Substituting~\eqref{eq:grad-q} into~\eqref{eq:grad-v-step1}:
\begin{align}
  \nabla v_\pi(s)
  &= \sum_a \nabla\pi(a\mid s)\,q_\pi(s,a) + \sum_a \pi(a\mid s)\,\gamma\sum_{s'}\mathcal{P}(s'\mid s,a)\,\nabla v_\pi(s') \notag \\
  &= \sum_a \nabla\pi(a\mid s)\,q_\pi(s,a) + \gamma\sum_{s'}\underbrace{\left[\sum_a \pi(a\mid s)\,\mathcal{P}(s'\mid s,a)\right]}_{=:\,P_\pi(s'\mid s)}\nabla v_\pi(s'). \label{eq:grad-v-substituted}
\end{align}
Define the shorthand
\[
  \phi(s) := \sum_a \nabla\pi(a\mid s)\,q_\pi(s,a),
\]
which collects all the ``direct gradient'' contributions at state~$s$, and let $P_\pi(s'\mid s) = \sum_a \pi(a\mid s)\,\mathcal{P}(s'\mid s,a)$ denote the one-step state-to-state transition probability under policy~$\pi$. Then~\eqref{eq:grad-v-substituted} becomes the recursion:
\begin{equation}\label{eq:grad-v-recursive}
  \nabla v_\pi(s) = \phi(s) + \gamma\sum_{s'\in\mathcal{S}} P_\pi(s'\mid s)\,\nabla v_\pi(s').
\end{equation}

\medskip\noindent\textbf{Step 4 (Unrolling the recursion).}\quad
Equation~\eqref{eq:grad-v-recursive} is a linear recursion in~$\nabla v_\pi$. We unroll it by repeatedly substituting on the right-hand side. After one substitution:
\begin{align*}
  \nabla v_\pi(s)
  &= \phi(s) + \gamma\sum_{s'} P_\pi(s'\mid s)\Bigl[\phi(s') + \gamma\sum_{s''} P_\pi(s''\mid s')\,\nabla v_\pi(s'')\Bigr] \\
  &= \phi(s) + \gamma\sum_{s'} P_\pi(s'\mid s)\,\phi(s') + \gamma^2\sum_{s''}\underbrace{\left[\sum_{s'} P_\pi(s'\mid s)\,P_\pi(s''\mid s')\right]}_{=\,P_\pi^{(2)}(s''\mid s)}\nabla v_\pi(s'').
\end{align*}
In general, writing $P_\pi^{(k)}(x\mid s) = \Pr\{S_{t+k}=x\mid S_t=s,\pi\}$ for the $k$-step transition probability, and continuing the expansion:
\begin{equation}\label{eq:unrolled}
  \nabla v_\pi(s) = \sum_{k=0}^{\infty}\gamma^k\sum_{x\in\mathcal{S}} P_\pi^{(k)}(x\mid s)\,\phi(x).
\end{equation}
Note that $P_\pi^{(0)}(x\mid s)=\mathbf{1}[x=s]$ (the identity), so the $k=0$ term recovers~$\phi(s)$ as expected. The series converges absolutely because $\|\phi(x)\|$ is bounded (it involves products of bounded quantities) and $\sum_k\gamma^k<\infty$.

\medskip\noindent\textbf{Step 5 (Applying to the performance measure).}\quad
The performance measure is $J(\bm{\theta})=v_\pi(s_0)$, so we set $s=s_0$ in~\eqref{eq:unrolled}:
\begin{align}
  \nabla J(\bm{\theta})
  &= \sum_{k=0}^{\infty}\gamma^k\sum_{x\in\mathcal{S}} P_\pi^{(k)}(x\mid s_0)\,\phi(x) \notag \\
  &= \sum_{x\in\mathcal{S}} \underbrace{\left[\sum_{k=0}^{\infty}\gamma^k P_\pi^{(k)}(x\mid s_0)\right]}_{=:\,\eta(x)} \phi(x) \notag \\
  &= \sum_s \eta(s)\sum_a q_\pi(s,a)\,\nabla\pi(a\mid s). \label{eq:pgt-eta-form}
\end{align}
The interchange of sums is justified by absolute convergence. Finally, since $\eta(s)\ge 0$ and $\sum_s\eta(s)>0$, we may normalise:
\begin{equation}\label{eq:pgt-final}
  \nabla J(\bm{\theta}) \propto \sum_{s\in\mathcal{S}} \mu_\pi(s)\sum_{a\in\mathcal{A}} q_\pi(s,a)\,\nabla_{\bm{\theta}}\,\pi(a\mid s,\bm{\theta}),
\end{equation}
where $\mu_\pi(s)=\eta(s)/\sum_{s'}\eta(s')$ is the on-policy distribution defined in~\eqref{eq:mu-pi}.
\end{proof}

\section{Discussion}
\label{sec:ch4-discussion}

\paragraph{Why the state distribution disappears.} The central insight of the proof is that $\nabla_{\bm{\theta}}\mu_\pi$ does not appear in the final expression. Intuitively, the effect of~$\bm{\theta}$ on the state distribution is absorbed into the weighting by~$\eta(s)$ through the recursive unrolling in Step~4. This absorption is possible precisely because $\mathcal{P}$ and~$\mathcal{R}$ are independent of~$\bm{\theta}$ (assumption~(iii)), so that the only $\bm{\theta}$-dependence in the Bellman equation for~$q_\pi$ enters through~$v_\pi$. Were the transition dynamics to depend on~$\bm{\theta}$---as they effectively do in the multi-agent setting, where the ``environment'' includes other agents' evolving policies---the step from~\eqref{eq:grad-q} would produce additional terms, and the clean recursive structure would break down. This is precisely the complication that the Meta-MAPG theorem of Chapter~\ref{ch:meta-mapg} addresses.

\paragraph{Connection to practical algorithms.} The Policy Gradient Theorem is what makes algorithms like \textsc{reinforce}~\cite{williams1992} and actor-critic methods work. The key practical consequence is that the gradient~\eqref{eq:pgt-final} can be written as an expectation under the agent's own behaviour: $\nabla J \propto \mathbb{E}_{s\sim\mu_\pi,\,a\sim\pi}[q_\pi(s,a)\,\nabla\ln\pi(a\mid s)]$ via the log-derivative trick (see Section~\ref{subsec:reinforce}). This expectation can be estimated from sample trajectories without ever computing $\mu_\pi$ explicitly. In practice, replacing $q_\pi(s,a)$ with the sample return~$G_t$ yields REINFORCE; replacing it with the TD error~$\delta_t$ yields actor-critic. The theorem guarantees that both are unbiased estimators of the true gradient (up to a positive scalar), despite having access only to sampled experience.

\paragraph{The proportionality constant.} The proportionality in~\eqref{eq:pgt-final} hides a positive normalisation factor $\sum_s\eta(s) = 1/(1-\gamma)$. In the episodic case with $\gamma=1$, this factor becomes the expected episode length and the proportionality becomes an equality when $\eta(s)$ is replaced by the (unnormalised) expected number of visits~\cite{sutton2018rl}. Since gradient ascent is invariant to positive rescaling of the gradient, the proportionality does not affect the optimisation.


\input{chapters/ch06-meta-mapg-theorem}

\chapter{Lean Formalisation}
\label{app:lean}

Both main results have been formalised in Lean~4 with Mathlib dependencies. Source files: \texttt{formalization/EvidenceWeightedPG.lean} (585 lines) and \texttt{formalization/OpponentShapingPG.lean} (467 lines).

\section{What Is Verified}

The formalisation axiomatises the stochastic game framework (policy space, projection, gradient signal decomposition) and \emph{verifies} the following steps:

\begin{enumerate}
    \item \textbf{Non-expansive projection}: $\|\operatorname{proj}(x) - \operatorname{proj}(y)\| \leq \|x - y\|$ (axiom, used in energy inequality).
    \item \textbf{Energy inequality} (Lemma~\ref{lem:energy-ch3}): expanding $D_{n+1}$ and separating drift, noise, bias, and gradient magnitude terms.
    \item \textbf{Variance scaling} (Lemma~\ref{lem:var-scaling-ch3}): block-diagonal structure of $W_n$ gives $\|W_n U_n\|^2 = \sum_i w_i^2 \|U_i\|^2$.
    \item \textbf{AM-HM inequality} (Theorem~\ref{thm:am-hm-ch3}): Cauchy--Schwarz applied to $(\sqrt{V_i})$ and $(1/\sqrt{V_i})$, yielding $\operatorname{HM}(V) \leq \operatorname{AM}(V)$.
    \item \textbf{Bias absorption} (Lemma~\ref{lem:bias-absorption-ch3}): triangle inequality for $\|b_n + \lambda_n \mathrm{OS}(\pi_n)\|$.
    \item \textbf{Spectral contraction}: if $S + \lambda S_H \preceq -\mu_{\mathrm{LOLA}} I$, then $\langle v^{\mathrm{LOLA}}(\pi), \pi - \pi^* \rangle \leq -\mu_{\mathrm{LOLA}} \|\pi - \pi^*\|^2$.
\end{enumerate}

\section{What Is Axiomatised}

The following are axiomatised (declared as \texttt{axiom} or left as \texttt{sorry}), as their full formalisation would require an extensive measure-theoretic library:
\begin{itemize}
    \item Martingale convergence (Doob's submartingale convergence theorem).
    \item Gladyshev's lemma (almost sure convergence of $D_n$).
    \item Chung's lemma (rate extraction from the Lyapunov recursion).
    \item The dominated convergence interchange $\nabla_\theta \int = \int \nabla_\theta$ in the REINFORCE derivation.
\end{itemize}

\section{Representative Code}

The core contraction step for EWPG (simplified):

\begin{verbatim}
-- Under SOS, weighted gradient provides contraction.
theorem ewpg_contraction
    (hSOS : forall pi in ball pi_star rho,
      inner (v pi) (pi - pi_star) <= -mu * norm(pi - pi_star)^2)
    (hw : forall i, w_min <= w i) :
    inner (W * v pi) (pi - pi_star) <=
      -mu * w_min * norm(pi - pi_star)^2 := by
  calc inner (W * v pi) (pi - pi_star)
      = sum i, w i * inner (v_i pi) (pi_i - pi_star_i) := ...
    _ <= w_min * sum i, inner (v_i pi) (pi_i - pi_star_i) := ...
    _ <= -mu * w_min * norm(pi - pi_star)^2 := ...
\end{verbatim}

The formalisation demonstrates that the key algebraic arguments are machine-checkable. Completing the measure-theoretic foundations would require Mathlib contributions for stochastic approximation theory, which we leave as future work.


\bibliographystyle{book}
\bibliography{myreferences}

\end{document}
